The splitting: $k \notin V_1$ since $k(u) = 1$, and $V_1$ has codimension one in $T^*M$, so $T^*M = V_1 \oplus \R k$ and the displayed decomposition of $T$ follows. Orthogonality: for $\alpha \in V_1$, $h(\alpha, \nO) = \nu\beta\,\alpha(u) = 0$ and likewise against $\nInf$; and $h(\alpha, k) = G(\alpha, k) = 0$ because $k$ is radical on $\Sig$. (1) $V_1$ meets the radical $\R k$ trivially, so the quotient map $V_1 \to T^*M/\shL$ is an isomorphism intertwining $G|_{V_1}$ with the induced form $\overline{G}$, which is nondegenerate of signature $(p{-}1,q)$ by Theorem \claimref{sheaf.degeneration-sheaf}(2). (2) The entries: $h(k,k) = G(k,k) = 0$; $h(k, \nO) = \nu \cdot k(u) = \nu$; $h(k, \nInf) = \mu$; the adjoined plane carries $H$. Expanding the determinant along the first row gives $-\nu(\,\nu \cdot 0 - (-1)\mu\,) + \mu(\,\nu \cdot (-1) - 0\,) = -\mu\nu - \mu\nu = -2\mu\nu$. (3) $k^2 = \nO^2 = 0$ and $h(k, \nO) = \nu \neq 0$ is a hyperbolic plane. Orthogonality of $w$: $h(w, k) = 0 + \nu\mu - \mu\nu = 0$ and $h(w, \nO) = \nu + \nu(-1) - 0 = 0$. Its square: $h(w,w) = k^2 + \nu^2 \nInf^2 + \mu^2 \nO^2 + 2\nu\,h(k,\nInf) - 2\mu\,h(k,\nO) - 2\mu\nu\,h(\nInf,\nO) = 2\nu\mu - 2\mu\nu + 2\mu\nu = 2\mu\nu$. A hyperbolic plane contributes $(1,1)$; the complementary line contributes the sign of $2\mu\nu$. (4) Add the signatures of the orthogonal summands (\claimref{std-algebra.sylvester-inertia}): $(p{-}1, q) + (1,1) + (\operatorname{sign} \mu\nu)$.
